\paragraph{4. Symmetry of Information from Average-Case Easiness of NP.}
In this section, we present a formal proof of Theorem 1.2. In fact, under a plausible assumption that $\mathsf E$ requires exponential-size circuits to compute, we prove that SoI follows from an approximation algorithm for time-bounded Kolmogorov complexity.
Theorem 4.1 immediately implies Theorem 1.2.
\begin{proof}[Proof of Theorem 1.2]
It is known that the two hypotheses of Theorem 4.1 are implied by the assumption that $\mathsf{NP}$ is easy on average: Buhrman et al. [18] showed that $\mathsf E\not\subseteq\bigcap_{\epsilon>0}\mathsf{i.o.SIZE}(2^{\epsilon n})$ if $\mathsf{DistNP}\subseteq\mathsf{AvgP}$; We also have $\mathsf{GapMINKT}\in\mathsf P$ if $\mathsf{DistNP}\subseteq\mathsf{AvgP}$ by Lemma 2.1.
\end{proof}
We now present the formal proof of Theorem 4.1.
\begin{proof}[Proof of Theorem 4.1]
Let $\widetilde K$ be the polynomial-time algorithm of Fact 3.4 which satisfies the property that there exists a polynomial $p$ such that for every $x\in\{0,1\}^*$ and every $t\geq|x|$,
\[K^{p(t)}(x)-\log p(t)\leq\widetilde K(x;1^t)\leq K^t(x).\tag{9}\]
Fix strings $x\in\{0,1\}^n$ and $y\in\{0,1\}^m$ and an integer $t\geq n+m$. The proof of SoI is given by analyzing the following three values for $z\sim\{0,1\}^{nk}$, $w\sim\{0,1\}^{nk+k}$, $z'\sim\{0,1\}^{m\ell}$, and $w'\sim\{0,1\}^{m\ell+\ell}$:
\[\begin{gathered}
\widetilde K(\mathsf{DP}_k(x;z),\mathsf{DP}_\ell(y;z');1^{t'}),\\
\widetilde K(w,\mathsf{DP}_\ell(y;z');1^{t'}),\\
\widetilde K(w,w';1^{t'}),
\end{gathered}\]
where $t'=t^{O(1)}$, $k\approx K^t(x,y)-\ell$, and $\ell\approx K^{\operatorname{poly}(t)}(y)$ are parameters chosen later.
First, observe that Fact 3.2 implies that $K(w,w')\geq|w|+|w'|-2$ with probability at least $3/4$. Let $\theta:=|w|+|w'|-2-\log p(t')$; using Equation (9), we obtain
\[\Pr_{w,w'}[\widetilde K(w,w';1^{t'})\geq\theta]\geq\frac34.\tag{10}\]
Next, we set the parameter $\ell$ to $K^{p'(t)}(y)-\log p'(t)-1$, where $p'$ is some large polynomial. Consider a randomized circuit $D$ that takes $w'$ as input as well as random bits $w$ and outputs 1 if and only if $\widetilde K(w,w';1^{t'})\geq\theta$. By the contrapositive of Lemma 2.2, $\mathsf{DP}_\ell(y;-)$ is a pseudorandom generator secure against $D$; i.e., $D$ cannot distinguish $\mathsf{DP}_\ell(y;z')$ and $w'$ in the sense that
\[\left|\Pr_{w,z'}[\widetilde K(w,\mathsf{DP}_\ell(y;z');1^{t'})\geq\theta]-\Pr_{w,w'}[\widetilde K(w,w';1^{t'})\geq\theta]\right|<\frac14,\]
which, together with Equation (10), implies that
\[\Pr_{w,z'}[\widetilde K(w,\mathsf{DP}_\ell(y;z');1^{t'})\geq\theta]\geq\frac12.\]
Finally, we compare $\widetilde K(w,\mathsf{DP}_\ell(y;z');1^{t'})$ with $\widetilde K(\mathsf{DP}_k(x;z),\mathsf{DP}_\ell(y;z');1^{t'})$. On one hand, since $|w|=|z|+k$ and $|w'|=|z'|+\ell$, we have
\[\Pr_{w,z'}[\widetilde K(w,\mathsf{DP}_\ell(y;z');1^{t'})\geq|z|+|z'|+k+\ell-2-\log p(t')]\geq\frac12.\tag{11}\]
On the other hand, observe that for some $t':=\operatorname{poly}(t)$,
\begin{align*}
\widetilde K(\mathsf{DP}_k(x;z),\mathsf{DP}_\ell(y;z');1^{t'})
&\leq K^{t'}(\mathsf{DP}_k(x;z),\mathsf{DP}_\ell(y;z'))\\
&\leq K^t(x,y)+|z|+|z'|+O(\log n)
\end{align*}
holds because the strings $\mathsf{DP}_k(x;z)$ and $\mathsf{DP}_\ell(y;z')$ can be computed from $k,\ell,z,z'$, and a program of size $K^t(x,y)$ that outputs $(x,y)$ in time $t$. We now set $k:=K^t(x,y)-\ell+O(\log t)$ so that
\[\Pr_{z,z'}[\widetilde K(\mathsf{DP}_k(x;z),\mathsf{DP}_\ell(y;z');1^{t'})<|z|+|z'|+k+\ell-2-\log p(t')]=1.\tag{12}\]
Let $D_y$ be a randomized circuit that takes an input $w$ and random bits $z'$ and outputs 1 if and only if $\widetilde K(w,\mathsf{DP}_\ell(y;z');1^{t'})<|z|+|z'|+k+\ell-2-\log p(t')$. It follows from Equations (11) and (12) that
\[\Pr_{z,z'}[D_y(\mathsf{DP}_k(x;z);z')=1]-\Pr_{w,z'}[D_y(w;z')=1]\geq1-\frac12=\frac12.\]
Using Lemma 2.2, we obtain
\[K^{\operatorname{poly}(t)}(x\mid D_y)\leq k+O(\log t)=K^t(x,y)-K^{p'(t)}(y)+O(\log t).\]
It follows that for some large polynomial $q$,
\begin{align*}
K^{q(t)}(x\mid y)&\leq K^t(x\mid D_y)+K^t(D_y\mid y)+O(1)\\
&\leq K^t(x,y)-K^{p'(t)}(y)+O(\log t)\\
&\leq K^t(x,y)-K^{q(t)}(y)+\log q(t)
\end{align*}
as desired.
\end{proof}
